% ================== Chapitre 5 : cycle de vie et exploitation ===============
\chapter{Cycle de vie des enregistrements et exploitation}
\label{ch:cycle}

\section{Propriétaire d'écriture unique et séquence de publication}
\label{sec:proprietaire}

Le document 03 pose une règle explicite : ne choisir qu'un propriétaire
d'écriture par enregistrement. Il ouvre deux pistes, la mise à jour dynamique
issue de Kea au moyen du RFC 2136, ou la mise à jour par l'agent au travers de
l'API autoritative.

\begin{table}[H]
\centering
\small
\caption{Arbitrage du mécanisme de mise à jour}\label{tab:arbitrage}
\begin{tabular}{@{}R{0.22}R{0.30}R{0.33}R{0.15}@{}}
\hautfilet
\textbf{Option} & \textbf{Apport} & \textbf{Limite} & \textbf{Décision}\\
\medfilet
Mise à jour par l'agent via l'API autoritative &
  Un seul émetteur pour toutes les zones. La modification porte sur un RRset
  complet, ce qui rend l'opération naturellement convergente. Le retour d'API
  fournit un état exploitable par la machine d'états et corrélable à un
  \opt{operation\_id}. &
  Ajoute une dépendance à la disponibilité de l'agent pour toute publication. &
  \textbf{Retenue}\\
Mise à jour dynamique RFC 2136 déclenchée par Kea &
  Rapproche la publication de l'attribution du bail et supprime un
  intermédiaire. &
  Introduit un second émetteur, donc un conflit possible de propriété sur les
  mêmes enregistrements. Exige une clé TSIG par zone, soit cent treize clés à
  gérer, et une compatibilité du moteur de stockage qui reste à vérifier. &
  Alternative documentée\\
\basfilet
\end{tabular}
\end{table}

\begin{figure}[H]
\centering
\figSequence
\caption{Séquence de publication d'un enregistrement et branche dégradée}
\label{fig:sequence}
\end{figure}

\label{sec:publication}
Les deux appels ci-dessous constituent le contrat que l'agent doit respecter.
Ils sont donnés pour une machine du tenant 7 recevant l'adresse
\ip{10.30.70.37}, qui appartient bien au pool dynamique \ip{.10} à \ip{.199}
défini par le document 02.

\begin{lstlisting}[caption={Publication conjointe de l'enregistrement direct et
  de l'enregistrement inverse},label=lst:publication]
PATCH /api/v1/servers/localhost/zones/tenant-007.gandal.internal.  HTTP/1.1
Host: ns1.gandal.internal:8443
X-API-Key: <cle injectee depuis le magasin de secrets>
Content-Type: application/json

{
  "rrsets": [{
    "name": "vm-0042.tenant-007.gandal.internal.",
    "type": "A",
    "ttl": 300,
    "changetype": "REPLACE",
    "records": [ { "content": "10.30.70.37", "disabled": false } ],
    "comments": [ { "account": "agent",
                    "content": "operation_id=7f3c..., lease_id=..." } ]
  }]
}

PATCH /api/v1/servers/localhost/zones/70.30.10.in-addr.arpa.  HTTP/1.1

{
  "rrsets": [{
    "name": "37.70.30.10.in-addr.arpa.",
    "type": "PTR",
    "ttl": 300,
    "changetype": "REPLACE",
    "records": [ { "content": "vm-0042.tenant-007.gandal.internal.",
                   "disabled": false } ]
  }]
}
\end{lstlisting}

La publication n'est pas considérée comme acquise à la réception du code 204.
L'agent interroge ensuite ns1, ns2, rec1 et rec2 pour le nom publié et ne fait
progresser l'opération vers l'état \opt{SERVICES\_READY} qu'après concordance
des quatre réponses. Ce contrôle détecte en particulier un transfert de zone
bloqué, défaut qui resterait sinon invisible jusqu'à la première panne de ns1.

\section{Idempotence, rejeu et réconciliation}
\label{sec:idempotence}

L'exigence \exig{net-nfr-006} demande un rejeu déterministe et idempotent. Le
choix d'un changement de type \opt{REPLACE} portant sur un RRset complet y
répond par construction : en notant $\Phi$ l'application d'une publication sur
l'état des zones, $\Phi \circ \Phi = \Phi$, c'est-à-dire que l'état final ne
dépend pas du nombre d'exécutions. Une précaution reste nécessaire.

\begin{encadre}
\titreencadre{Précaution sur le rejeu}
Une modification par l'API incrémente le numéro de série de la zone, y compris
lorsque le contenu soumis est identique à celui déjà publié. Un rejeu
systématique provoquerait donc autant de transferts de zone inutiles vers ns2.

\smallskip
L'agent doit par conséquent lire le RRset avant de l'écrire et n'émettre la
modification que si le contenu diffère. Le rejeu reste correct dans tous les
cas, mais il devient également économe.
\end{encadre}

Une réconciliation périodique compare les allocations enregistrées dans l'IPAM
aux enregistrements effectivement publiés, dans les deux sens. Elle signale les
divergences sans les corriger automatiquement : une correction automatique sur
un état mal compris est le moyen le plus rapide de propager une erreur à
l'ensemble des zones. Deux familles de divergence sont attendues,
l'enregistrement publié sans allocation correspondante, trace d'une suppression
incomplète, et l'allocation active sans enregistrement, trace d'une publication
interrompue.

\section{Suppression, quarantaine et mode dégradé}
\label{sec:suppression}

L'ordre de suppression est l'inverse de l'ordre de publication et il n'est pas
interchangeable. Le document 03 impose d'attendre le retrait des interfaces
virtuelles et la fin des références avant libération, et précise qu'une panne
DNS ne libère pas une adresse encore utilisée.

\begin{enumerate}
  \item Retrait de l'interface virtuelle de la machine et fin du bail.
  \item Suppression des enregistrements A et PTR par un changement de type
    \opt{DELETE} sur les deux zones concernées.
  \item Attente de la durée de quarantaine $Q = 900$ secondes, qui satisfait
    l'inégalité \eqref{eq:ttl}.
  \item Retour de l'adresse au pool par l'IPAM.
\end{enumerate}

Si la suppression des enregistrements échoue, l'adresse n'est pas rendue au pool
et l'opération reste dans l'état \opt{DELETING}, observable. Le comportement
symétrique s'applique à la publication : une panne de l'API ou des serveurs
autoritatifs laisse l'opération en cours, signalée par une alerte, sans que
l'adresse déjà attribuée soit considérée comme libre. Le trafic des machines
déjà configurées n'est pas interrompu : seules les opérations nouvelles sont
suspendues. Cette distinction est précisément ce que le test \opt{G04} doit
établir.

\section{Amorçage, ordre de reprise et adresses de secours}
\label{sec:amorcage}

Le document Underlay insiste sur la rupture des cycles de démarrage : si le
serveur DNS et le bastion dépendent uniquement d'un ensemble dont l'accès
suppose ces mêmes services, une panne générale rend la restauration difficile.
Le socle minimal ci-dessous doit donc être présent sur les trois hôtes, sur XO,
sur le bastion et sur les machines de services.

\begin{lstlisting}[language=pdnsconf,caption={Adresses de secours à conserver
  hors du service DNS},label=lst:hosts]
# /etc/hosts  : socle de reprise GANDAL
# Permet d'administrer et de reconstruire l'infrastructure sans dependre
# du service DNS lui-meme. Non publie : voir export-etc-hosts=no.

127.0.0.1    localhost
10.30.0.11   h1.infra.gandal.internal         h1
10.30.0.12   h2.infra.gandal.internal         h2
10.30.0.13   h3.infra.gandal.internal         h3
10.30.0.20   xo.infra.gandal.internal         xo
10.30.0.21   bastion.infra.gandal.internal    bastion
10.30.4.10   ns1.gandal.internal              ns1
10.30.4.11   ns2.gandal.internal              ns2
10.30.4.12   rec1.gandal.internal             rec1
10.30.4.13   rec2.gandal.internal             rec2
10.30.4.20   kea1.infra.gandal.internal       kea1
10.30.4.21   kea2.infra.gandal.internal       kea2
10.30.4.30   ntp1.infra.gandal.internal       ntp1
10.30.4.31   ntp2.infra.gandal.internal       ntp2
10.30.4.40   ipam.infra.gandal.internal       ipam
10.30.4.41   agent.infra.gandal.internal      agent
\end{lstlisting}

\begin{longtable}{@{}C{0.09}R{0.24}R{0.15}R{0.52}@{}}
\caption{Ordre de reprise après indisponibilité étendue}\label{tab:reprise}\\
\hautfilet
\textbf{Rang} & \textbf{Élément remis en service} & \textbf{Dépend de} &
\textbf{Remarque}\\
\medfilet
\endfirsthead
\hautfilet
\textbf{Rang} & \textbf{Élément remis en service} & \textbf{Dépend de} &
\textbf{Remarque}\\
\medfilet
\endhead
\basfilet
\endfoot
1 & Transport et hôtes & Aucune &
  Les hyperviseurs se joignent par adresse de gestion, sans résolution.\\
2 & Source de temps & Transport &
  La validation DNSSEC et la vérification des certificats exigent une horloge
  plausible. Au démarrage à froid, le client de temps doit être autorisé à
  corriger un écart important avant toute vérification cryptographique.\\
3 & Autoritatifs ns1 puis ns2 & Temps &
  Le fichier \opt{hosts} suffit à les atteindre.\\
4 & Résolveurs rec1 et rec2 & Autoritatifs &
  Leur propre résolution locale pointe sur la boucle locale, jamais sur leur
  homologue, afin d'éviter une dépendance croisée.\\
5 & IPAM puis agent & DNS &
  La réconciliation ne démarre qu'une fois le DNS joignable.\\
6 & DHCP & DNS et IPAM &
  Les options distribuées référencent les résolveurs.\\
7 & Orchestration et supervision & Ensemble des précédents &
  Dernière étape ; la supervision constate l'état sans conditionner la remise
  en service.\\
\end{longtable}

Le point sensible est le rang 2. La validation cryptographique du temps par NTS
suppose un certificat valide, dont la vérification suppose une horloge correcte,
qui suppose elle-même le service de temps. Le document 03 signale ce cas sans le
trancher. Du point de vue du DNS, la conséquence est simple : si l'horloge est
fausse, la validation DNSSEC rejette des réponses publiques pourtant légitimes,
ce qui se manifeste par des codes \opt{SERVFAIL} sur l'espace public alors que
l'espace interne répond normalement. Ce symptôme doit figurer dans la procédure
de diagnostic, car il oriente immédiatement vers l'horloge et non vers le DNS.

\section{Journalisation, métriques et alertes}
\label{sec:journaux}

La journalisation est placée sur les résolveurs, seuls points où l'adresse
source est celle du client réel. Elle est horodatée en temps universel
coordonné, conformément à \exig{net-fr-018}, et limitée aux éléments
nécessaires : horodatage, adresse source, nom demandé, type, code de réponse et
étiquette de vue. La durée de conservation et la responsabilité des données
relèvent de l'interface avec IAM et les responsables métier.

\begin{longtable}{@{}R{0.30}R{0.30}R{0.40}@{}}
\caption{Indicateurs et alertes du service DNS}\label{tab:alertes}\\
\hautfilet
\textbf{Indicateur} & \textbf{Seuil ou condition} & \textbf{Portée}\\
\medfilet
\endfirsthead
\hautfilet
\textbf{Indicateur} & \textbf{Seuil ou condition} & \textbf{Portée}\\
\medfilet
\endhead
\basfilet
\endfoot
Disponibilité de chaque instance & Absence de réponse à une sonde SOA pendant
  60 s & Majeure\\
Écart de numéro de série entre ns1 et ns2 & Écart persistant au-delà de deux
  cycles de transfert, soit 120 s &
  Majeure : la réplication est rompue et la redondance n'est plus effective\\
Taux de réponses \opt{REFUSED} par vue & Hausse significative par rapport à la
  référence &
  Mineure : politique trop stricte, mauvaise distribution des résolveurs, ou
  tentative d'accès\\
Taux de réponses \opt{SERVFAIL} & Hausse sur l'espace public &
  Majeure : dérive d'horloge, perte de la sortie externe ou échec de
  validation\\
Âge du dernier transfert réussi sur ns2 & Supérieur à deux fois
  \opt{refresh}, soit 7\,200 s & Majeure\\
Opérations de publication en échec & Toute opération bloquée au-delà de son
  délai & Majeure\\
Divergence IPAM et zones publiées & Toute divergence persistant après deux
  cycles de réconciliation & Mineure, à instruire\\
Tentatives de transfert de zone refusées & Toute occurrence &
  À instruire en sécurité\\
Absence de données de collecte & Interruption de la lecture en pull &
  Majeure : la panne de collecte doit alerter, et non effacer silencieusement
  les traces\\
\end{longtable}

\section{Sauvegarde, restauration et gestion des secrets}
\label{sec:sauvegarde}

L'objectif historique de reconstruction en moins de soixante minutes reste à
évaluer, et cette évaluation relève de T009. Ce dossier en fournit les éléments :
ce qui doit être sauvegardé, où, et dans quel ordre la restauration s'effectue.

\begin{lstlisting}[caption={Export et restauration de l'état autoritatif},
  label=lst:sauvegarde]
# --- export quotidien de l'etat autoritatif, depuis ns1
install -d -m 0750 /var/backups/gandal/dns
for Z in $(pdnsutil list-all-zones) ; do
  pdnsutil list-zone "$Z" > "/var/backups/gandal/dns/${Z}.zone"
done

# --- l'export est verse dans le depot versionne, sans aucun secret
#     Les cles TSIG et les cles d'API restent dans le magasin de secrets.

# --- restauration de ns1 sur une machine neuve
pdnsutil import-tsig-key gandal-xfr hmac-sha256 "$TSIG_GANDAL_XFR"
for F in /var/backups/gandal/dns/*.zone ; do
  Z=$(basename "$F" .zone)
  pdnsutil create-zone "$Z" ns1.gandal.internal
  pdnsutil load-zone   "$Z" "$F"
  pdnsutil set-meta    "$Z" TSIG-ALLOW-AXFR gandal-xfr
  pdnsutil set-meta    "$Z" SOA-EDIT-API    DEFAULT
done
pdnsutil rectify-all-zones
pdnsutil check-all-zones
\end{lstlisting}

\begin{table}[H]
\centering
\small
\caption{Gestion des secrets et des sauvegardes du service}\label{tab:secrets}
\begin{tabular}{@{}R{0.22}R{0.28}R{0.50}@{}}
\hautfilet
\textbf{Élément} & \textbf{Emplacement} & \textbf{Règle}\\
\medfilet
Clé TSIG de transfert & Magasin de secrets &
  Rotation annuelle ou sur incident. La rotation se fait en déclarant la
  nouvelle clé sur les deux serveurs avant de retirer l'ancienne.\\
Clé d'API de ns1 & Magasin de secrets, injectée dans l'environnement du
  service &
  Rotation lors de tout changement du porteur de l'agent. L'empreinte peut
  figurer en configuration à partir de la version 4.7.\\
Certificat TLS du mandataire & PKI interne, interface IF-EXT-01 &
  Renouvellement avant échéance ; la supervision du délai restant fait partie
  des indicateurs attendus.\\
Exports de zones & Dépôt versionné et copie hors du cluster &
  Une copie qui ne serait disponible que dans le cluster administré n'offre
  aucune reprise autonome après sa perte.\\
\basfilet
\end{tabular}
\end{table}
